%% LLM-SoccerArena -- coauthor bullet draft for KDD 2026.
%% Every bullet is intended to contain final-paper-level content. The remaining
%% work is prose conversion, final-data refresh, and resolution of explicit TODOs.
\documentclass[sigconf,review]{acmart}

\let\Bbbk\relax
\usepackage{amssymb}
\usepackage{array}
\usepackage{enumitem}
\usepackage{subcaption}
\usepackage{tabularx}
\usepackage{xspace}
\usepackage{tikz}
\usetikzlibrary{positioning,arrows.meta,shapes.geometric,calc,fit,backgrounds}
\usepackage{pgfplots}
\pgfplotsset{compat=1.18}

\newcommand{\arena}{\mbox{LLM-SoccerArena}\xspace}
\newcommand{\arenalink}{\href{https://llm-soccerarena.com}{llm-soccerarena.com}}
\newcommand{\NumModels}{8}

\definecolor{draftblue}{RGB}{34,91,145}
\definecolor{draftlight}{RGB}{239,246,252}
\definecolor{todoorange}{RGB}{176,91,0}
\definecolor{todolight}{RGB}{255,246,230}

\newenvironment{draftbullets}
  {\begin{itemize}[leftmargin=1.25em,itemsep=1.5pt,topsep=2pt,parsep=0pt,partopsep=0pt]}
  {\end{itemize}}
\newcommand{\TODO}[1]{{\color{red}#1}}
% \newcommand{\TODO}[1]{%
%   \par\smallskip\noindent
%   \fcolorbox{todoorange}{todolight}{%
%     \parbox{\dimexpr\linewidth-2\fboxsep-2\fboxrule}{\small\textbf{TODO:} #1}}%
%   \par\smallskip}

\newcommand{\takeaway}[1]{%
  \par\smallskip\noindent
  \fcolorbox{draftblue}{draftlight}{%
    \parbox{\dimexpr\linewidth-2\fboxsep-2\fboxrule}{\small\textbf{Takeaway.} #1}}%
  \par\smallskip}

\makeatletter
\newcommand{\daggerauthornote}[1]{%
  \authornotemark[2]%
  \g@addto@macro\@authornotes{\footnotetext[2]{#1}}%
}
\makeatother

\newcommand{\draftstatus}{%
  \fcolorbox{todoorange}{todolight}{%
    \parbox{\dimexpr\linewidth-2\fboxsep-2\fboxrule}{\small\textbf{Coauthor bullet draft (July 2026).}
    The argument, paragraph contents, figures, and provisional findings are included at
    final-paper level. Exact prose, the final 104-match refresh, two unresolved tournament
    questions, rationale adjudication, and explicitly marked citation or presentation decisions
    remain open. \textbf{TODO:} Once the result artifacts are frozen, convert the bullets to
    prose and compress the main text and floats to the eight-page submission limit.}}}

% The bullet draft contains deliberately information-dense lines and long compound
% labels. Give TeX a little extra flexibility without changing the ACM type size.
\setlength{\emergencystretch}{2em}

\setcopyright{acmlicensed}
\copyrightyear{2026}
\acmYear{2026}
\acmDOI{XXXXXXX.XXXXXXX}
\acmConference[KDD '27]{Proceedings of the 33rd ACM SIGKDD Conference on Knowledge Discovery and Data Mining}{August 2027}{San Jose, USA}
\acmISBN{978-1-4503-XXXX-X/2026/08}

\begin{document}

\title{\arena: Benchmarking LLMs on Real-World Predictions in Sports}

\author{Jonas Schr\"oder}
\daggerauthornote{Joint first authors.}
\affiliation{\institution{LMU Munich} \city{Munich}\country{Germany}}
\affiliation{\institution{MCML}\city{Munich}\country{Germany}}
\email{jonas.schroeder@lmu.de}

\author{Jonas Schweisthal}
\authornotemark[2]
\affiliation{\institution{LMU Munich}\city{Munich}\country{Germany}}
\affiliation{\institution{MCML}\city{Munich}\country{Germany}}
\email{jonas.schweisthal@lmu.de}

\author{Oliver M\"uller}
\affiliation{\institution{Paderborn University}\city{Paderborn}\country{Germany}}
\email{oliver.mueller@uni-paderborn.de}

\author{Markus Weinmann}
\affiliation{\institution{University of Cologne}\city{Cologne}\country{Germany}}
\email{weinmann@wiso.uni-koeln.de}

\author{Stefan Feuerriegel}
\affiliation{\institution{LMU Munich}\city{Munich}\country{Germany}}
\affiliation{\institution{MCML}\city{Munich}\country{Germany}}
\email{feuerriegel@lmu.de}

\renewcommand{\shortauthors}{Schr\"oder and Schweisthal, et al.}

\begin{abstract}
Large language models (LLMs) are increasingly used to reason about and support decisions concerning
uncertain future events, yet this ability is difficult to evaluate. Static benchmarks risk having
answers already present in training data, are often constructed only after outcomes are known, and
rarely record what information a model could access when making a prediction. We introduce
\href{https://llm-soccerarena.com/}{\arena}, an open-source, real-time benchmark and public platform
that records timestamped LLM forecasts before outcomes occur. We demonstrate \arena in the context
of the 2026 FIFA World Cup. Across all 104 matches, \arena recorded 9,984 match-level forecast
attempts (9,316 valid) from eight contemporary proprietary and open-weight LLMs, including GPT-5.5,
Gemini 3.1 Pro, and Claude Opus 4.8. It additionally collected 480 forecasts for 15 tournament-level
questions, such as the group winners, semifinalists, top-scorer team, and eventual champion. We vary
four conditions: (1) \emph{model family}; (2) \emph{information access}, i.e., whether the model may
search the web or must rely on internal knowledge alone; (3) \emph{prompting style}, i.e., whether it
is asked for a scoreline first or for home/draw/away probabilities first; and (4) \emph{forecast
timing}, from the start of a tournament stage to 24 and 2 hours before kickoff. Forecasts are
evaluated using proper scoring rules and accuracy-based metrics, alongside measures of output
validity, tool use, latency, tokens, and cost. We find that: (1) web access reliably improves
probability forecasts, lowering the Brier score by 0.023 (95\% CI: 0.004--0.040), while no pairwise
model difference within the complete seven-model panel remains significant after correcting for
multiple comparisons; (2) prompting style does not improve overall forecast quality, but asking for
probabilities first produces approximately four percentage points more draw scorelines; (3) models
assign highly similar probabilities to the same matches (mean pairwise correlation $r = 0.943$), so
averaging their forecasts yields only a small additional improvement; and (4) higher self-reported
confidence is associated with more accurate subsequent predictions, although it is not itself a
probability of being correct. \arena therefore evaluates not only which LLM forecasts best, but how
information access, prompting, and timing shape real-world LLM forecasting systems.
\end{abstract}

\begin{CCSXML}
<ccs2012>
 <concept><concept_id>10010147.10010257</concept_id><concept_desc>Computing methodologies~Machine learning</concept_desc><concept_significance>500</concept_significance></concept>
 <concept><concept_id>10010147.10010178.10010179</concept_id><concept_desc>Computing methodologies~Natural language processing</concept_desc><concept_significance>300</concept_significance></concept>
</ccs2012>
\end{CCSXML}
\ccsdesc[500]{Computing methodologies~Machine learning}
\ccsdesc[300]{Computing methodologies~Natural language processing}
\keywords{large language models, forecasting, live benchmark, information access, sports analytics, uncertainty}

\maketitle
\draftstatus

\section{Introduction}

\begin{draftbullets}
  \item Forecasting platforms turn uncertain future events into an observable and evaluable process.
  They define questions and resolution criteria, collect timestamped probability judgments, update
  them as information changes, and score them once outcomes become known. Their public websites add an
  important communication layer: forecasts become inspectable before resolution, leaderboards make
  performance comparable, and historical records allow users to follow both individual predictions
  and progress over time. For LLMs, this infrastructure is especially valuable because it evaluates
  an economically and societally relevant capability---supporting analysis, planning, and decisions
  under uncertainty---through commitments about the real world rather than answers to a fixed test.

  \item Recent live benchmarks demonstrate how these functions can be implemented at scale.
  ForecastBench maintains a changing pool of questions about future events, evaluates probabilistic
  forecasts from machine and human forecasters, and communicates their performance through public
  leaderboards \citep{forecastbench}. TS-Arena applies a forecast pre-registration protocol to live
  time-series streams: predictions are submitted before the corresponding observations exist and are
  evaluated only when the future becomes the past \citep{tsarena}. Together, such platforms provide
  continuously renewable test material, enforce temporal ordering between prediction and outcome,
  and make model progress visible outside a one-off paper evaluation.

  \item This prospective design addresses fundamental weaknesses of conventional static benchmarks.
  When questions and answers are assembled from historical material, evaluation targets or closely
  related reports may already occur in model training data. Even if direct memorization is absent,
  running a benchmark after an event has resolved permits hindsight leakage through retrieved pages,
  latent model knowledge, or benchmark construction choices. Static datasets also rarely preserve
  the exact forecast origin: when the prediction was made, which evidence was then available, which
  model version produced it, and whether external tools were used. Committing forecasts while targets
  remain genuinely unknown makes the temporal boundary explicit and substantially strengthens the
  validity of real-world forecasting evaluation \citep{tsarena,forecastbench}.

  \item Yet prospective collection alone does not identify why one LLM forecaster outperforms
  another. Existing forecasting research asks whether language models can produce useful real-world
  probabilities \citep{halawi_etal_2024}, but an operational forecaster is more than a model name. Web
  retrieval, search behavior, prompt order, structured-output constraints, sampling settings, and the
  time remaining until resolution can all affect the prediction. Open submission leaderboards are
  useful for comparing the best complete systems, including scaffolds and tools, but their scores may
  consequently combine model capability with unequal information access and implementation choices.
  Conversely, tightly closed model tests may improve comparability while omitting precisely the
  deployment conditions that determine practical performance. What is missing is a design that varies
  these conditions systematically while holding the forecasting task and evaluation pipeline fixed.

  \item A second gap concerns comparability across questions and over time. Broad collections of
  heterogeneous real-world events provide valuable coverage, but differences in topic, horizon,
  outcome frequency, resolution rules, and available public information can obscure whether a score
  difference reflects forecasting skill or a different mix of questions. Moreover, a single forecast
  per event cannot show whether a model benefits from late-breaking evidence or whether rankings are
  stable across information horizons. A homogeneous sequence of events with common targets, repeated
  forecast times, and identical experimental conditions complements broad benchmarks by enabling
  within-event and within-model comparisons rather than relying only on aggregate leaderboard ranks.

  \item A third gap lies between a public benchmark website and an auditable research artifact.
  Public interfaces necessarily prioritize concise forecasts, summary scores, rankings, and timely
  updates. Scientific analysis additionally requires the records behind those displays: exact prompts
  and model versions, raw and validated responses, failed or repaired outputs, tool-execution traces,
  timestamps, latency, token use, and cost. These traces matter because apparently valid forecasts can
  still contain incoherent probability vectors or contradictions between a predicted category and an
  exact scoreline. Without a shared provenance layer, readers cannot determine whether a leaderboard
  difference originates in probabilistic judgment, elicitation, information retrieval, or output
  handling, and reported results become difficult to reproduce after providers update their models.

  \item We introduce \arena (\arenalink), a prospective benchmark and public website designed around
  these requirements and deployed for the 2026 FIFA World Cup. Before each match, the system records
  forecasts from multiple contemporary LLMs under a controlled $2\times2$ design: closed versus open
  book, crossed with score-first versus probabilities-first elicitation. It repeats collection at the
  opening of each tournament stage, T--24h, and T--2h, retaining the actual lead time of every call.
  The match sequence supplies a common three-way probability target and exact scoreline target, while
  15 prospectively elicited tournament-level questions extend evaluation to longer-horizon outcomes.

  \item The platform connects scientific control with public inspectability. A shared SQLite archive
  stores prompts, raw and validated outputs, model/version identifiers, timing, web-search traces,
  latency, token use, and cost. The website exposes pre-match forecasts, conditions, and live
  leaderboards from this archive; after collection, the research pipeline freezes and verifies the
  same records before computing reported results. The public platform is therefore not a separate
  demonstration but a transparent view onto the benchmark's underlying evidence.

  \item The paper contributes: (i) a deployed, timestamped infrastructure for contamination-resistant
  real-world forecasting; (ii) a factorial benchmark that separates model choice from information
  access and elicitation order; (iii) repeated forecast horizons and a homogeneous event domain for
  controlled within-event comparisons; (iv) a multi-metric evaluation of probability, categorical,
  and scoreline quality alongside validity, tool use, latency, and cost; (v) evidence on retrieval
  benefits, prompting side effects, forecast similarity, and structured probability failures; and
  (vi) a reproducible release protocol linking every reported number and public display to the frozen
  prediction archive.
\end{draftbullets}

\TODO{After the final freeze, update the abstract and contribution bullets with 104 matches and all 15 resolved tournament questions. Optional: add one recent citation on tool-augmented LLM forecasting after extending references.bib.}

\section{Related Work}

\begin{draftbullets}
  \item \textbf{Dynamic evaluation and contamination resistance.} Static benchmarks deteriorate as
  their test items enter training corpora, model developers optimize against known leaderboards, and
  later systems gain access to information unavailable when the benchmark was created. Dynamic NLP
  evaluation has responded with human-and-model-in-the-loop data creation in Dynabench, hosted
  evaluation and operational metrics in Dynaboard, and regularly refreshed objective questions in
  LiveBench \citep{dynabench,dynaboard,livebench}. Recent work systematizes this shift from static to
  dynamic benchmarking and emphasizes temporal validity, renewable items, and transparent update
  protocols \citep{chen_etal_2025_contamination}. ForecastBench transfers these principles to
  probabilistic event forecasting by continuously evaluating unresolved future questions and
  comparing machine, public, and expert forecasts \citep{forecastbench}. TS-Arena instead
  pre-registers numerical forecasts on live energy streams before target observations exist
  \citep{tsarena}. ExAnte shows why this distinction matters: even explicit historical cutoffs do not
  reliably prevent LLMs from drawing on internalized post-cutoff knowledge \citep{exante}.

  \item \textbf{LLMs as real-world event forecasters.} Autocast established a large benchmark of
  forecasting-tournament questions paired with a date-indexed news corpus, allowing historical
  decision conditions to be reconstructed \citep{autocast}. Autocast++ improves this setting through
  recency-aware retrieval and ranking \citep{autocastpp}, while retrieval, reasoning, fine-tuning, and
  ensembling can bring end-to-end LLM systems closer to competitive human aggregates
  \citep{halawi_etal_2024}. Other studies compare prompting and forecasting strategies with human
  predictions \citep{pratt_etal_2024}, equip agents with search and numerical tools
  \citep{hsieh_etal_2024}, or evaluate tool-using agents on international relations and long-context
  news events \citep{mirai,tcelongbench}. More recent benchmarks broaden the design space: Prophet
  Arena continuously collects live tasks and decomposes the forecasting pipeline
  \citep{prophetarena}; WorldReasoner evaluates outcomes, time-valid evidence, and causal reasoning
  jointly \citep{worldreasoner}; and PolyBench couples timestamp-locked prediction-market, news, and
  order-book states with trading-oriented metrics \citep{polybench}. These studies establish
  real-world forecasting as a distinct LLM capability, but differ substantially in temporal protocol,
  task heterogeneity, permitted information, and whether they compare base models or optimized agents.

  \item \textbf{Retrieval, agents, and observable tool use.} Retrieval-augmented generation separates
  parametric model knowledge from an updateable external memory \citep{lewis_etal_2020}. ReAct
  interleaves reasoning with actions against external environments \citep{yao_etal_2023}, and
  Toolformer studies when and how a model learns to invoke APIs \citep{schick_etal_2023}. General
  tool-agent benchmarks further show that successful completion depends on tool selection, argument
  construction, execution, and recovery from errors, not simply on nominal tool availability
  \citep{wang_etal_2024_gta}. This distinction is central to forecasting: enabling web search is an
  experimental condition, whereas actual search execution, retrieved context, latency, tokens, and
  cost are operational outcomes. We therefore treat open versus closed book as a controlled factor
  while separately measuring provider-observable tool compliance and resource use.

  \item \textbf{Probabilistic evaluation, calibration, and confidence.} Proper scoring rules reward
  both accuracy and honest uncertainty. The Brier score and ranked probability score provide classic
  foundations for binary, multiclass, and ordered outcomes \citep{brier_1950,epstein_1969}, with
  decomposition results separating reliability, resolution, and uncertainty \citep{murphy_1973}.
  General treatments formalize proper scoring rules and the complementary roles of calibration and
  sharpness \citep{gneiting_raftery_2007,gneiting_etal_2007}. Neural networks can be systematically
  miscalibrated even when their classification accuracy is high \citep{guo_etal_2017}; for LLMs,
  self-evaluation and verbalized confidence can contain useful information but depend on task,
  elicitation method, and model alignment \citep{kadavath_etal_2022,tian_etal_2023,xiong_etal_2024}.
  Accordingly, \arena distinguishes proper-score performance from reliability, categorical and exact
  score decisions, self-reported confidence, and structural coherence of probability outputs.

  \item \textbf{Prompt sensitivity and structured elicitation.} LLM behavior can change with prompt
  format, label priors, and example ordering even when the underlying task is unchanged
  \citep{zhao_etal_2021,lu_etal_2022}. Forecasting adds a further elicitation problem because a system
  may be asked for a probability distribution, a most likely category, an exact scoreline, and a
  rationale in one structured response. The order of these requests may anchor later fields or expose
  disagreements between outcome-level and score-level modes. Our score-first versus
  probabilities-first intervention therefore leaves definitions and the JSON schema unchanged while
  reversing only the requested reasoning order, and evaluates distributional side effects alongside
  average predictive performance.

  \item \textbf{Forecast diversity and aggregation.} Combining forecasts often improves robustness,
  although the benefit depends on error dependence and information overlap \citep{clemen_1989}.
  Methods for aggregating probabilities include linear and nonlinear opinion pools, logit-based
  extremization, and distributional combinations designed to manage calibration and dispersion
  \citep{gneiting_ranjan_2013,satopaa_etal_2014,lichtendahl_etal_2013}. Human forecasting tournaments
  similarly show that training, teamwork, aggregation, and cognitive strategy can improve geopolitical
  judgment \citep{tetlock_2005,mellers_etal_2014}. Multiple LLMs, however, may share training data,
  retrieved sources, and inductive biases, so a larger model panel does not necessarily provide
  independent information. We therefore measure same-event probability similarity and test
  same-condition ensembles rather than assuming a mechanical ``wisdom of the model crowd.''

  \item \textbf{Football prediction.} Football offers a mature set of specialized baselines.
  Classical approaches model goals through independent or dependent Poisson processes
  \citep{maher_1982,dixon_coles_1997,karlis_ntzoufras_2003,koopman_lit_2015}; ratings and regression
  models incorporate current team strength, time depreciation, and tournament covariates
  \citep{hvattum_arntzen_2010,ley_etal_2019,groll_etal_2015}; and random forests and broader machine
  learning frameworks use richer match and contextual features
  \citep{schauberger_groll_2018,sportsprediction,beal_etal_2021}. Bookmaker-derived consensus models
  provide strong external probability baselines for tournaments
  \citep{leitner_etal_2010,zeileis_etal_2018}. \arena does not propose another football-specific model.
  It uses a common sequence of objectively resolved matches to compare general-purpose LLM systems
  under prospectively recorded information and elicitation conditions.

  \item \textbf{Research gap.} Table~\ref{tab:benchmark_comparison} summarizes the closest benchmark
  designs. To our knowledge, no prior platform combines prospective collection on unresolved events,
  a homogeneous repeated event domain, a crossed model--information-access--elicitation design,
  multiple operational forecast horizons, and joint evaluation of probabilities, categories, exact
  scorelines, and set-valued tournament questions. Existing work also rarely links the public
  leaderboard to a frozen archive containing raw prompts and responses, validation outcomes, search
  traces, model versions, latency, tokens, and cost. \arena fills this gap by treating the deployed
  forecasting configuration---rather than the model name alone---as the object of comparison while
  retaining the controls and provenance needed to attribute performance differences.
\end{draftbullets}

\begin{table*}[t]
  \centering
  \caption{Comparison with representative forecasting and dynamic-evaluation benchmarks. ``Public
  layer'' distinguishes a live or browsable benchmark interface from a released dataset or codebase.}
  \label{tab:benchmark_comparison}
  \scriptsize
  \setlength{\tabcolsep}{3pt}
  \renewcommand{\arraystretch}{1.16}
  \begin{tabularx}{\textwidth}{@{}p{2.15cm}p{2.45cm}p{2.75cm}p{2.85cm}p{1.65cm}X@{}}
    \hline
    \textbf{Benchmark} & \textbf{Forecast target} & \textbf{Temporal protocol} &
    \textbf{Information / controlled factors} & \textbf{Public layer} &
    \textbf{Primary distinction from \arena} \\
    \hline
    Autocast \citep{autocast} & Diverse event probabilities and numerical outcomes &
    Historical simulation using date-indexed news & Fixed corpus; retrieval can be studied &
    Dataset and code & Resolved static archive; no live collection or factorial deployment design \\
    ForecastBench \citep{forecastbench} & Broad binary future-event probabilities &
    Dynamic questions unresolved at forecast time & Base-model and open system submissions; human
    comparisons & Live leaderboard & Heterogeneous question pool; no common scoreline target or
    crossed access-by-elicitation design \\
    TS-Arena \citep{tsarena} & Numerical energy time series & Live forecast pre-registration &
    Standardized model service; no LLM search intervention & Live portal & Time-series foundation
    models rather than judgmental LLM event forecasts \\
    MIRAI \citep{mirai} & International-relation event types & Historical cutoff over resolved events &
    Structured database and news APIs; agent strategies & Dataset, code, and demo & Tool-use benchmark
    on relational events; not prospectively collected \\
    ExAnte \citep{exante} & Stock, QA, event, and publication tasks & Simulated ex-ante cutoff &
    Measures latent use of post-cutoff knowledge & Dataset and code & Leakage diagnosis rather than a
    deployed live forecasting platform \\
    Prophet Arena \citep{prophetarena} & Live real-world event probabilities and market returns &
    Continuous live-task collection & Modular knowledge, context, and reasoning stages & Public site &
    Broad market-oriented tasks rather than one repeated domain with scoreline outputs \\
    PolyBench \citep{polybench} & Binary prediction markets and trading returns & Timestamp-locked
    market snapshots & Synchronized news and order-book state & Dataset and code & Trading and market
    microstructure focus; no prompt-order or closed/open-book factorial \\
    \textbf{\arena} & \textbf{H/D/A probabilities, exact scores, and tournament sets} &
    \textbf{Prospective stage-opening, T--24h, and T--2h collection} & \textbf{Model $\times$ access
    $\times$ elicitation, with observed tool traces} & \textbf{Live website and leaderboard} &
    \textbf{Homogeneous sports domain linked to a frozen, auditable prediction archive} \\
    \hline
  \end{tabularx}
\end{table*}

\section{Why Football? Why the 2026 World Cup?}

\begin{draftbullets}
  \item Managerial and strategic decisions frequently depend on forecasts made under incomplete,
  changing, and sometimes conflicting information. Football provides a bounded real-world decision
  environment in which those central difficulties remain visible: forecasters must combine prior
  knowledge with new evidence, quantify uncertainty, and commit to a prediction before the relevant
  outcome exists. This makes the domain useful for studying forecast-informed decision support rather
  than merely measuring sports knowledge.

  \item At the same time, football offers unusually clean evaluation. Fixtures and kickoff times are
  known, 90-minute outcomes are objectively observed, and the three-way home/draw/away target can be
  scored immediately with proper rules. The same forecast yields interpretable structured outputs
  such as scorelines, goal totals, and knockout advancement, allowing multiple aspects of a system's
  decision support to be assessed against one realized event.

  \item The task remains genuinely challenging. Team form, injuries, lineups, travel, tactics,
  rankings, and market beliefs evolve before kickoff; signals can be incomplete or contradictory;
  draws and low-scoring matches preserve substantial irreducible uncertainty. These properties make
  information access an operationally meaningful benchmark condition and create strong external
  anchors---statistical models, bookmakers, and expert systems---for interpreting forecast quality.

  \item The expanded 2026 World Cup supplies 104 matches across group and knockout stages, 48 teams
  with heterogeneous information availability, and a fixed structure supporting prospective questions
  about group winners, semifinalists, the champion, and the top-scoring team. It is therefore a
  complex but controlled case study: broad enough to expose model and deployment differences, yet
  bounded enough to audit every forecast and outcome. The benchmark does not claim that football
  represents every management setting; it provides a transparent test bed for mechanisms that recur
  in decisions under uncertainty.
\end{draftbullets}

\TODO{Optional extension (Oli/Markus): connect the decision-making motivation more explicitly to management and decision-support research and strengthen the football-domain rationale. Suitable references should be added to references.bib before citation. Add an official FIFA citation for the 2026 format, 104-match schedule, and competition structure.}

\section{Platform, Website, and Benchmark}
\label{sec:benchmark}

\begin{draftbullets}
  \item \textbf{One archive, two interfaces.} \arena is both a prospective forecast-collection system
  and a public interface over the same prediction archive. The website displays forecasts and live
  evaluations; the paper pipeline analyzes frozen records from the database. Public viewing and
  engagement are not experimental units and do not enter any model comparison.

  \item \textbf{Controlled design.} The stable prediction key is match $\times$ model/version $\times$
  horizon $\times$ access $\times$ prompt $\times$ sample. Closed-book calls disable tools; open-book
  calls enable search and explicitly request current information, while actual search execution is
  logged rather than assumed. Score-first and probabilities-first prompts require the same output
  schema but reverse the requested elicitation order. Stage opening, T--24h, and T--2h are operational
  snapshots with their actual lead times retained.

  \item \textbf{Prospective data flow.} Fixtures enter the prompt builder before provider execution,
  schema validation, deterministic normalization where permitted, and at most one repair attempt.
  SQLite stores raw prompts and responses, hashes, model versions, timing, validity, tool traces,
  tokens, latency, and cost. Official 90-minute outcomes enter only after forecast storage; the
  analysis begins from a consistent database backup and derives typed Parquet tables linked by a
  hash manifest.

  \item \textbf{Public inspectability.} The match view exposes H/D/A probabilities, predicted
  scorelines, rationales, conditions, and timing. The analytics view provides a multi-metric
  leaderboard, filters, and cumulative performance across the calendar, while tournament views expose
  prospective questions and the bracket. The site is an interface to the benchmark, not a separate
  analytical data source.
\end{draftbullets}

\begin{figure*}[ht]
  \centering
  \begin{subfigure}[t]{0.98\textwidth}
    \centering
    \resizebox{0.96\linewidth}{!}{\input{generated/fig_flow}}
    \caption{Prospective collection and post-outcome evaluation share one auditable archive.}
  \end{subfigure}

  \vspace{0.5em}
  \begin{subfigure}[t]{0.47\textwidth}
    \centering
    \includegraphics[width=\linewidth]{figures/screenshot-card.png}
    \caption{Match forecasts and their deployment conditions.}
  \end{subfigure}\hfill
  \begin{subfigure}[t]{0.47\textwidth}
    \centering
    \includegraphics[width=\linewidth]{figures/screenshot-analytics.png}
    \caption{Filterable multi-metric analytics and trajectories.}
  \end{subfigure}
  \caption{\arena architecture and public interface. Forecasts cross the outcome barrier only after
  they have been timestamped and stored; the website and reproducible analysis read the same
  underlying records.}
  \label{fig:platform}
\end{figure*}

\TODO{Refresh the two website crops after the final interface update and replace the current flow diagram with the final combined architecture artwork if needed. Add final model/version strings and release URL/DOI.}

\section{Evaluation Protocol}
\label{sec:protocol}

\begin{draftbullets}
  \item The universal primary target is the official 90-minute home, draw, or away result, and T--24h is
  the primary operational snapshot. The seven-model complete panel is used for common comparisons;
  invalid forecasts are not imputed. Fable is reported as an operational-reliability result with a
  common-match sensitivity panel, not as a main leaderboard entry. The 15 tournament questions remain
  question-level prospective units.

  \item Main performance reporting includes the unscaled three-class Brier score (0 best, 2 worst,
  uniform forecast $2/3$), log loss in nats, probability-modal H/D/A accuracy with fractional credit
  for tied maxima, exact-score accuracy, and Kicktipp points. Paired contrasts first average
  prespecified cells within each match and then estimate the mean paired difference.

  \item Uncertainty uses 10,000 group/knockout-stratified match-bootstrap replicates for 95\%
  bootstrap-$t$ intervals and 10,000 sign-flip permutations on the same mean-difference estimand,
  with master seed 20260715. Declared contrast families use Holm adjustment. Timing sensitivities,
  leave-one-match-out checks, multiplicity families, and acceptance gates are detailed in the
  appendix.

  \item The SQLite database is frozen, integrity-checked, and hashed before analysis. The public CSV
  is used once for reconciliation, never as an analytical input. External odds, Opta, and classical
  model candidates enter only if their provenance, feature timing, target, and reproduction gates
  pass; the current Results section does not depend on them.
\end{draftbullets}

\TODO{Replace all provisional counts after the final 104-match freeze. Complete and attach the final acceptance checklist before converting result bullets into claims.}

\section{Results}

\subsection{Overall Performance}

\begin{draftbullets}
  \item The leaderboard evaluates the complete panel at T--24h, averages the four access-by-prompt
  cells equally within each model--match, and then averages matches. This makes the match the analysis
  unit and prevents cells with greater coverage from receiving additional weight. All displayed
  intervals resample matches within group/knockout strata.

  \item No model leads every metric. Gemini currently has the lowest Brier score (0.500) and log loss
  (0.844), DeepSeek the highest probability-modal H/D/A accuracy (65.5\%), and GPT the highest
  exact-score accuracy (16.8\%) and mean Kicktipp score (1.47). The Brier range is narrow
  (0.500--0.536), so the benchmark yields a metric-dependent ordering rather than a universal winner.

  \item Only Gemini versus Qwen currently survives Holm adjustment across the 21 paired Brier model
  comparisons. Most between-model probability-quality differences are therefore small relative to
  match-to-match uncertainty even though descriptive ranks differ.

  \item Cumulative Brier and modal accuracy are volatile early in the tournament and stabilize as
  matches accumulate. Match-level median Brier currently rises from 0.403 in the group stage to 0.591
  in the quarterfinals, but late-stage samples are small; the stage panel describes realized forecast
  error rather than establishing a monotonic difficulty trend.

  \item Fable produced 552 valid forecasts from 1,220 planned calls, with 668 API errors. Its
  incomplete coverage excludes it from the common leaderboard; the 26-match common-set analysis is
  retained only as an operationally selected sensitivity result.
\end{draftbullets}

\begin{table*}[t]
  \centering
  \caption{Complete-panel T--24h performance over 101 completed matches. Values are model means with
  95\% match-bootstrap intervals; lower Brier and log loss are better, while higher values are better
  for the remaining metrics. Ranks are descriptive and will be refreshed after the final freeze.}
  \label{tab:leaderboard}
  \resizebox{\textwidth}{!}{\input{generated/overall_leaderboard}}
\end{table*}

\begin{figure*}[t]
  \centering
  \includegraphics[width=\textwidth]{figures/analysis/fig_overall_performance.pdf}
  \caption{Overall performance across the tournament. Cumulative trajectories use the same balanced
  model--match aggregation as Table~\ref{tab:leaderboard}; stage distributions show one complete-panel
  mean per match with raw observations and stage sample sizes.}
  \label{fig:overall}
\end{figure*}

\takeaway{Model rankings depend on the metric: probability quality, categorical hits, and scoreline performance favor different systems, while most paired Brier differences remain uncertain after familywise adjustment.}

\subsection{RQ1: Information Access}

\begin{draftbullets}
  \item Open-book access improves average T--24h probability quality. Mean Brier is currently 0.5309
  closed book and 0.5054 open book, giving a paired closed-minus-open difference of 0.0255
  [0.0067, 0.0430], Holm-adjusted $p=0.025$. This corresponds descriptively to a 4.8\% reduction from
  the closed-book mean.

  \item The effect remains stable when any one match is omitted (leave-one-match-out range
  0.0235--0.0280) and when only calls made within $1{,}440\pm90$ minutes of kickoff are retained
  (0.0244 [0.0048, 0.0426], $n=98$). It is especially pronounced in the round of 32 (0.0611);
  the chronological view shows a rise into early knockouts and subsequent decline, with substantial
  local variation.

  \item Operational snapshots move in the expected direction but remain uncertain. Open-book Brier
  improves by 0.0064 from stage opening to T--24h and by another 0.0020 from T--24h to T--2h; both
  intervals include zero. Closed-book forecasts change very little across the same snapshots, which
  is the expected design-validity pattern when no new information enters. Actual calls center at
  24.33 and 2.10 hours before kickoff.

  \item Search availability does not imply uniform use. Observed search ranges from 48.7\% for Gemini
  and 61.0\% for GPT to at least 98.4\% for DeepSeek, Grok, Mistral, and Qwen. Open book adds on
  average 22,383 input tokens, 884 output tokens, 3.92 seconds, and USD~0.110 per forecast; model-level
  added cost ranges from roughly USD~0.014 to USD~0.321.

  \item Access gains also differ across models, from approximately zero for Qwen to 0.038 for GPT.
  The descriptive variance decomposition attributes 49\% of cell-mean variation to condition,
  37\% to model, and 14\% to their interaction, with broad bootstrap distributions. The evidence
  supports treating information access as a first-class benchmark dimension, not claiming that it
  dominates all between-model differences.
\end{draftbullets}

\begin{figure*}[t]
  \centering
  \includegraphics[width=\textwidth]{figures/analysis/fig_rq1_access.pdf}
  \caption{Information access. (A) Overall and stage-specific closed-minus-open T--24h Brier effects;
  positive values favor open book. (B) Match-level effects over the tournament calendar with a
  descriptive smoother. (C) Changes across operational snapshots and the closed-book negative
  control. (D) Observed search use and incremental cost by model.}
  \label{fig:rq1}
\end{figure*}

\takeaway{Open-book access improves T--24h probabilistic forecasts on average and survives timing and match-deletion checks, but gains vary and require additional search, latency, tokens, and cost.}

\subsection{RQ2: Elicitation Strategy}

\begin{draftbullets}
  \item Score-first prompting requests a scoreline before probabilities and expected goals;
  probabilities-first prompting reverses that requested order while retaining the same definitions
  and JSON schema. Comparisons match horizon and access within each match. Displayed advantages are
  oriented so that positive values favor probabilities first.

  \item Elicitation order has essentially no average Brier advantage: probabilities first minus the
  score-first baseline yields an advantage of $-0.0008$ [$-0.0045$, 0.0031], Holm-adjusted
  $p=0.679$. No registered advantage appears for log loss, modal H/D/A accuracy, exact-score
  accuracy, or Kicktipp points. Modal accuracy directionally favors score first by 0.92 percentage
  points, but this does not survive the four-metric Holm family.

  \item The design shows a descriptive crossover. Score first is directionally better closed book,
  while probabilities first is directionally better open book. The additional probabilities-first
  advantage under open book is 0.0074 [0.0007, 0.0146], with raw $p=0.041$ and Holm-adjusted
  $p=0.083$.

  \item Forecast structure nevertheless changes with prompt order. The probabilities-first outputs contain
  3.82 percentage points more score-implied draws and 3.89 points more exact 1--1 scorelines, while
  assigning 0.20 points less mean draw probability. A draw score paired with a non-draw modal outcome
  consequently rises by 5.16 points. Six of seven complete-panel models show the increased
  score-draw direction.

  \item This is not necessarily a logical inconsistency: one exact scoreline can be modal even when
  the aggregate draw outcome is not. Score/outcome modal agreement is 70.2\% versus 64.1\% for
  closed-book score-first versus probabilities-first calls, and 84.7\% versus 80.5\% open book.
  Repairs remain rare and no probability normalization is applied.
\end{draftbullets}

\begin{figure*}[t]
  \centering
  \includegraphics[width=\textwidth]{figures/analysis/fig_rq2_elicitation.pdf}
  \caption{Elicitation effects. The left panel reports registered performance contrasts; the right
  panel shows outcome-independent changes in score-implied draws, exact 1--1 predictions, mean draw
  probability, and draw-score/non-draw-modal-outcome divergence. Positive values favor
  probabilities-first elicitation.}
  \label{fig:rq2}
\end{figure*}

\takeaway{Prompt order barely changes average forecast quality but changes forecast structure: probabilities-first produces more 1--1 scorelines while slightly lowering draw probability.}

\subsection{RQ3: Calibration of Elicited Probabilities}

\begin{draftbullets}
  \item Reliability is evaluated from one complete-panel probability vector per
  match--access--prompt cell: currently 404 vectors over 101 completed matches. Five
  approximately equal-frequency bins are shown separately for home, draw, and away outcomes, with
  pointwise intervals obtained by resampling whole matches.

  \item Home reliability varies around the diagonal without a simple directional pattern. Draw
  forecasts occupy a narrower range (14.9\%--30.9\%). Mean draw probability is 24.5\% against an
  observed draw rate of 27.7\%; the predicted-minus-observed gap is $-3.2$ percentage points
  [$-12.7$, 5.1], $p=0.477$, which is compatible with mild underprediction but does not establish a
  systematic aggregate gap.

  \item The clearest localized departures occur for low and middle away probabilities: predicted
  bin means of 8.3\%, 17.5\%, and 25.2\% correspond to observed frequencies of 0.0\%, 4.9\%, and
  6.2\%. The highest away bin reverses descriptively (61.5\% predicted versus 77.8\% observed), a
  pattern consistent with probabilities being pulled toward the center.

  \item The self-reported confidence field contains a useful exploratory relative-difficulty signal.
  After ranking confidence within each model--access--prompt cell while preserving ties, mean Brier
  declines from 0.611 to 0.419 and modal H/D/A accuracy rises from 53.8\% to 73.6\% between the
  lowest and highest groups. These values describe association with subsequent accuracy; confidence
  is not interpreted as a calibrated probability of correctness.
\end{draftbullets}

\begin{figure*}[t]
  \centering
  \includegraphics[width=\textwidth]{figures/analysis/fig_rq3_calibration.pdf}
  \caption{Probability reliability and exploratory confidence association. Home, draw, and away
  probabilities use separate equal-scale reliability panels with pointwise match-bootstrap
  intervals. Confidence groups are defined within model and condition; lower Brier and higher modal
  accuracy indicate better forecasts.}
  \label{fig:rq3}
\end{figure*}

\takeaway{Probability reliability differs by outcome. Higher within-cell confidence accompanies better forecasts, but it is a relative difficulty signal rather than a probability of correctness.}

\subsection{RQ4: Forecast Diversity and Herding}

\begin{draftbullets}
  \item Similarity is measured only within the same T--24h match--access--prompt cell. Pairwise
  Jensen--Shannon (JS) divergence captures absolute separation between H/D/A probability vectors;
  class-specific correlations combined on the Fisher-$z$ scale capture whether models rank matches
  similarly. High correlation and non-zero JS can therefore coexist.

  \item The seven models share a strong forecast signal. Mean pairwise JS divergence is 0.00449 nats
  (range 0.00163--0.00854), while mean combined correlation is 0.943 (range 0.904--0.984). Claude
  Opus--GPT has the lowest JS divergence, Gemini--GPT the highest correlation, and several of the
  most distinct pairings include Mistral.

  \item Open-book access does not detectably increase convergence. Mean pairwise JS is 0.00461 closed
  book and 0.00438 open book; the paired open-minus-closed difference is
  $-0.2\times10^{-3}$ [$-1.0$, 0.6]$\times10^{-3}$ nats, $p=0.572$. We therefore describe
  forecast similarity rather than inferring behavioral copying or herding.

  \item Equal-weight ensembles are formed only inside one horizon--access--prompt cell. Ensemble
  Brier is 0.00478 lower than average member Brier by the Brier ambiguity identity, a small
  mechanical diversity gain. The ensemble directionally beats four models and trails three, but no
  ensemble-versus-model comparison survives the seven-test Holm family. Averaging improves
  robustness relative to the average member without reliably dominating individual systems.
\end{draftbullets}

\begin{figure*}[t]
  \centering
  \includegraphics[width=\textwidth]{figures/analysis/fig_rq4_diversity_herding.pdf}
  \caption{Forecast similarity and same-cell ensembles. Lower-triangle matrices report pairwise JS
  divergence and Fisher-$z$-combined probability correlation. The forest plot compares the
  equal-weight ensemble with the average member and each prespecified model without mixing horizons,
  access conditions, or prompts.}
  \label{fig:rq4}
\end{figure*}

\takeaway{Models share a strong forecast signal. Same-cell averaging provides a small robustness gain over the average member, but does not reliably outperform individual models.}

\subsection{RQ5: Reasoning and Tool Use}

\begin{draftbullets}
  \item This section studies observable provider traces and generated rationale text, not private
  chain of thought. In the current full corpus, 76.6\% of 4,880 open-book forecasts show observable
  search, 16.6\% explicitly show no search, and 6.8\% have unknown status; the unknown records occur
  only in partial Fable runs. Complete-panel search observability is 84.5\%. Closed-book search is
  not applicable because tools are disabled by design.

  \item Every non-empty rationale is coded by two independent automated systems: the released keyword
  lexicon and a fixed-prompt LLM annotator. All disagreements plus a fixed-seed sample of 50 agreements
  are exported without outcomes, scores, match identity, or model identity for blinded human
  adjudication. Final prevalence uses the resolved labels; automated agreement and Cohen's $\kappa$
  are reported as measurement validation.

  \item The analysis covers 9,092 non-empty rationales. All 8,540 complete-panel forecasts contain a
  rationale, whereas Fable contains 552 of 1,220 (45.2\%), balanced across access. Evidence categories
  describe what the generated text mentions; they do not establish that a source was consulted or
  that the evidence is factually correct.

  \item Rationale length already shows an operational contrast: the median is 28 words closed book
  and 57 words open book across available rationales (28 and 55 words in the complete panel). The
  open-book median is higher for every model, although both baseline length and the increase vary
  substantially. Length is treated as a surface property, not a proxy for reasoning quality.

  \item Two or three short matched excerpts will illustrate how a closed- and open-book rationale for
  the same model, match, horizon, and prompt differ in mentioned evidence. Examples are selected
  without forecast correctness and carry no quantitative claim.

  \item As a lightweight exploratory view, one t-SNE map will project fixed text embeddings of the
  rationale corpus into two dimensions, with model represented by color and open/closed access by
  marker shape. The map is descriptive and is used only to show whether broad model- or access-level
  organization is visually apparent.
\end{draftbullets}

\begin{figure*}[t]
  \centering
  \fcolorbox{todoorange}{todolight}{%
    \parbox[c][4.2cm][c]{\dimexpr\textwidth-2\fboxsep-2\fboxrule}{\centering
      \textbf{RQ5 figure placeholder}\par\smallskip
      Panel A: adjudicated evidence-category prevalence for closed versus open book, with direct
      percentage-point differences.\par\smallskip
      Panel B: one exploratory t-SNE map of fixed rationale embeddings; model = color,
      access = marker shape.\par\smallskip
      \textbf{TODO:} configure the annotation model and API key, run the automated export, complete
      blinded human adjudication, generate the resolved evidence panel, and create/cache the fixed
      embeddings and t-SNE coordinates.}}
  \caption{Planned reasoning and rationale-content display. The prevalence panel will use adjudicated
  labels over the full non-empty corpus; the t-SNE panel is an exploratory visualization without an
  inferential claim.}
  \label{fig:rq5}
\end{figure*}

\TODO{Complete the mandatory annotation checkpoint; insert final category rates, agreement statistics, two to three illustrative excerpts, the t-SNE artifact, and the final takeaway sentence. Candidate wording: observable search and rationale content vary across systems, while open-book rationales are longer and mention [resolved evidence categories] more often.}

\takeaway{Open result: objective traces already show substantial differences in tool compliance and rationale length; the evidence-content comparison remains gated until blinded adjudication is complete.}

\begin{figure*}[t]
  \centering
  \includegraphics[width=\textwidth]{figures/analysis/fig_rq6_tournament_questions.pdf}
  \caption{Prospective tournament questions. (A) Forecast distributions for the 12 group winners,
  with equal-weight question means, realized winners, correct-pick counts, and the 25\% uniform
  reference. (B) Exact semifinal-set recovery and mean team recovery over 32 forecasts.}
  \label{fig:rq6}
\end{figure*}

\subsection{RQ6: Tournament-Level Questions}

\begin{draftbullets}
  \item Testing long-horizon forecast capabilities of the LLMs. The analysis units are the 15 prespecified questions: 12 group winners, the four
  semifinalists, the world champion, and the team providing the tournament's top scorer. The 32
  model--access--prompt forecasts per question describe within-question variation and are not treated
  as independent observations. Thirteen questions are currently resolved; champion and top-scoring
  team remain unscored.

  \item Group forecasts are usually strong. Across the 12 group questions, equal-weight means are
  91.15\% correct final picks, Brier 0.2825, log loss 0.5498, and 60.39\% probability on the
  realized winner. For context, a uniform four-team forecast gives 25\%, Brier 0.75, and log loss
  1.386. Ten groups have 32 correct picks of 32, and Mexico has 30 of 32.

  \item Group K is a concentrated prospective failure: Colombia receives a mean realized-winner
  probability of 27.19\% (range 12\%--37\%), while all 32 final picks select Portugal. The benchmark
  thus captures favorite structure well overall but can reproduce a shared pre-tournament belief
  across every model and condition strongly enough to miss an upset.

  \item Fourteen of 32 forecasts recover the exact four-team semifinal set (43.75\%). The average
  forecast contains 3.375 of the four teams, or 84.38\% team recovery; two forecasts contain two,
  16 contain three, and 14 contain all four realized semifinalists.

  \item The supplied semifinal marginal probabilities reveal a separate format failure. Their median
  sum is the required 4.000, but the mean is 4.857 and the range is 1.555--11.700. The deviation is
  strongly model-specific. Exact-set accuracy is unaffected, while marginal Brier must be read
  alongside this coherence audit; the supplied probabilities are not retrospectively normalized.
\end{draftbullets}

\takeaway{Prospective tournament forecasts are usually strong but share a Group K miss. Semifinal recovery is high, yet supplied inclusion probabilities often violate the sum-to-four constraint.}

\TODO{After the tournament, score the champion and top-scoring-team questions, update all 15-question summaries, and replace provisional wording in the RQ6 takeaway.}

\section{Discussion}

\begin{draftbullets}
  \item \textbf{Forecasting systems, not model names alone.} The metric-dependent leaderboard and the
  access effect show why a model identifier is an incomplete description of a deployed forecaster.
  Model, condition, and their interaction all contribute to observed variation. Deployment effects
  can be comparable in magnitude to between-model variation even though the current data do not
  support a claim that conditions dominate every model difference.

  \item \textbf{Information has a measurable benefit and a measurable price.} The open-book T--24h
  benefit survives timing and match-deletion checks, and the directional snapshot pattern is
  consistent with changing information becoming useful nearer kickoff. Yet tool availability does
  not create a standardized intervention: providers differ in whether they search, how much context
  they retrieve, and the latency, token, and financial cost they incur.

  \item \textbf{Elicitation changes structure before performance.} Reversing score and probability
  elicitation barely changes average quality, but it changes draw scorelines, exact 1--1 predictions,
  and the relation between score and outcome modes. Benchmarking prompt interventions therefore
  requires semantic and distributional audits alongside aggregate performance metrics.

  \item \textbf{Decisions, probabilities, and confidence are distinct outputs.} Reliability differs by
  outcome, self-reported confidence tracks relative difficulty without being a calibrated correctness
  probability, and strong tournament picks coexist with a shared Group K miss and invalid sums of
  semifinal marginals. Structured forecast systems should validate probability vectors, set
  marginals, final picks, and confidence fields separately.

  \item \textbf{A correlated crowd offers limited robustness.} Models rank matches similarly while
  retaining small probability-level differences. Shared web access does not detectably reduce this
  diversity. Same-cell averaging mechanically improves on the average member but does not reliably
  beat individual models, making it a modest robustness tool rather than a new forecasting oracle.

  \item \textbf{Prospective collection exposes operational outcomes.} Fable's missing calls, incomplete
  search compliance, validation records, latency, cost, and rationale coverage are scientifically
  meaningful properties that a retrospective leaderboard would hide. Live benchmarks therefore need
  stable prediction keys, frozen snapshots, raw traces, explicit missingness policies, and an
  intent-to-treat view of tool-enabled calls.

  \item \textbf{Scope and limitations.} This is one tournament with 104 matches at final freeze, one
  sample per cell, provider-specific retrieval, incomplete Fable coverage, sparse late stages, and
  only 15 tournament questions. Rationale labels describe observable text rather than hidden
  reasoning or factuality. No external candidate currently passes the provenance and reproduction
  gate, so we make no comparative claim against betting markets, Opta, or classical models. Repeated
  prospective deployments are needed before generalizing beyond football; the benchmark evaluates
  forecasts and does not provide wagering advice.
\end{draftbullets}

\TODO{Insert the final RQ5 evidence-category synthesis after adjudication. Reassess the conditions-versus-models sentence mechanically after the final freeze: use the locked fallback sentence only if its two predicates become true.}

\section{Conclusion}

\begin{draftbullets}
  \item \arena combines a public forecasting website with a prospective, auditable benchmark that
  records not just predictions but the information conditions, elicitation choices, tool behavior,
  and operational costs that produced them.
  \item The current results show that retrieval access improves average probability quality,
  prompting can reshape outputs without improving aggregate scores, model forecasts remain highly
  correlated, and structured predictions can succeed categorically while failing probability-format
  checks. The resulting benchmark reframes real-world LLM forecasting as an evaluation of complete
  systems under time and information constraints.
\end{draftbullets}

\TODO{Rewrite the conclusion into one compact final paragraph after all results and the release statement are frozen.}

\section*{Acknowledgments}

OpenAI Codex and Anthropic Claude were used as assistive tools for software development, analysis
implementation, visualization, and manuscript drafting and editing; the authors conceived the study,
directed and verified the work, and take responsibility for the final content.

\TODO{Before submission, verify whether the ACM disclosure should be supplemented with exact tool and model versions, representative prompts, and post-generation editing details in an appendix or supplementary artifact.}

\bibliographystyle{ACM-Reference-Format}
\bibliography{references}

\clearpage
\appendix

\section{Detailed Evaluation and Reproducibility Protocol}

\begin{draftbullets}
  \item Freeze SQLite through its backup mechanism inside a read transaction; run
  \texttt{PRAGMA integrity\_check}; record UTC freeze time, SHA-256, schema version, table counts,
  source path, and repository commit. Final-paper mode requires 104 reconciled fixtures and one
  official 90-minute outcome per completed match.
  \item Derive typed Parquet tables with explicit schemas and UTC timestamps. Every table records the
  source database hash, configuration hash, code commit, row count, key uniqueness, and file hash.
  The public website CSV is read once for reconciliation; discrepancies are logged rather than
  resolved by silently preferring either source.
  \item For every paired contrast, construct one paired difference per match after prespecified
  aggregation. Resample whole matches within group and knockout strata for 10,000 studentized
  bootstrap replicates and run 10,000 sign-flip permutations on the same mean. Report the estimate,
  95\% interval, raw and Holm-adjusted $p$ where applicable, descriptive median, and sample sizes.
  \item Declared Holm families cover access/prompt/interactions, complete-panel model pairs,
  operational snapshots, eligible stages, ensemble-versus-model contrasts, and any passing external
  baseline family. No Wilcoxon test or formal equivalence margin is used.
  \item Every figure and table emits a companion JSON with its estimand, estimate, interval,
  $p$-values, counts, aggregation rule, configuration hash, source-table hashes, output hash, and
  manifest key. Package, Python, operating-system, R-runtime, and external-tool versions are recorded.
\end{draftbullets}

\TODO{Insert the final manifest summary, acceptance checklist, exact Holm-family table, timing distributions, and data-release DOI.}

\section{Supporting Results}

\subsection{Fable operational reliability}
\begin{draftbullets}
  \item Report planned calls, valid forecasts, API errors, repair counts, rationale coverage, latency,
  and cost. Keep Fable outside the main rank and label the 26-match common-set comparison as a
  sensitivity result rather than a replacement leaderboard.
\end{draftbullets}

\subsection{Model-specific information-access effects}
\begin{figure*}[h]
  \centering
  \includegraphics[width=\textwidth]{figures/analysis/fig_rq1_access_by_model.pdf}
  \caption{Descriptive model-specific distributions of T--24h closed-minus-open Brier differences.
  The pooled access contrast is prespecified; model-specific panels are descriptive extensions.}
\end{figure*}

\subsection{Prompt modal agreement}
\begin{figure}[h]
  \centering
  \includegraphics[width=\columnwidth]{figures/analysis/fig_rq2_coherence.pdf}
  \caption{Score/outcome modal agreement by access condition and elicitation strategy. Exact-score
  and aggregate-outcome modes need not coincide, so the display is not labeled a logical-coherence
  test.}
\end{figure}

\subsection{Rationale length and embedding visualization}
\begin{figure*}[h]
  \centering
  \includegraphics[width=\textwidth]{figures/analysis/fig_rq5_rationale_length_appendix.pdf}
  \caption{Model-specific median and interquartile range of rationale length by access condition.
  Fable's partial rationale coverage is reported explicitly.}
\end{figure*}

\TODO{If the main RQ5 figure becomes too dense, move the single exploratory t-SNE panel here without adding additional robustness panels or inferential claims.}

\subsection{Semifinal probability-sum audit}
\begin{figure*}[h]
  \centering
  \includegraphics[width=\textwidth]{figures/analysis/fig_rq6_semifinal_probability_sums_appendix.pdf}
  \caption{Supplied semifinal inclusion-probability sums by model and condition. The required sum is
  four; raw marginals are audited without retrospective normalization.}
\end{figure*}

\section{Rationale Annotation Checkpoint}
\begin{draftbullets}
  \item The keyword coder uses the released category lexicon. The LLM coder uses a configured exact
  model string, fixed prompt, temperature zero, and no access to outcomes, performance, or keyword
  labels. Raw provider responses are cached and frozen because temperature zero does not guarantee
  bit-identical regeneration.
  \item All complete-vector disagreements and 50 fixed-seed agreements are shuffled into
  \texttt{disagreements.csv}. The file contains only a blinded annotation ID, rationale text, both
  automated label vectors, and empty human-label fields. The pipeline halts until every human field
  is valid and the immutable-column hash and row identities match.
  \item Final rates use human labels for adjudicated rows and shared automated labels for the remaining
  agreements. Category prevalence, raw agreement, and Cohen's $\kappa$ are reported separately.
\end{draftbullets}

\TODO{Supply the annotation model/version and API-key environment configuration, run the checkpoint, complete blinded adjudication, and freeze the resolved corpus.}

\section{External-Baseline Gate}
\begin{draftbullets}
  \item Betting odds, Opta, and the five Markus model candidates remain excluded until source IDs,
  retrieval timing, raw snapshot hashes, orientation, pre-kickoff feature timing, target compatibility,
  redistribution terms, and reproduction are established.
  \item Bookmaker overround must be removed within bookmaker before median aggregation. Closing odds
  can be compared only with open-book T--2h forecasts. Frozen Markus forecasts, if reproduced without
  tournament fitting or updating, are mechanically matched to stage opening rather than T--24h.
  \item The bundled Markus predictions are never analytical inputs. The supplied R code must be run in
  its isolated pinned environment, and every reproduced candidate retained without post-hoc selection.
\end{draftbullets}

\TODO{Keep the external-baseline audit open. If no candidate passes, retain only the gate table and exclusion reasons in the appendix; do not consume main-paper Results space.}

\section{Prompt Components}

\subsection*{Closed-book header}
\begin{quote}\small
Access condition: CLOSED\_BOOK. Do not use internet search, browsing, tools, APIs, or external data
sources. Use only your internal model knowledge, general football reasoning, typical football score
distributions, and the fixture information provided below.
\end{quote}

\subsection*{Open-book header}
\begin{quote}\small
Access condition: OPEN\_BOOK. Before making the prediction, use the available web-search tool to
check current public information about this match and both teams. Base the final prediction on the
retrieved public information plus your football reasoning.
\end{quote}

\subsection*{Score-first instruction}
\begin{quote}\small
First predict the most likely scoreline for the match. Then provide probabilities, expected goals,
and full-match/advancement probabilities that are consistent with that predicted scoreline.
\end{quote}

\subsection*{Probabilities-first instruction}
\begin{quote}\small
First estimate calibrated probabilities for the 90-minute result and expected goals. Then derive the
most likely scoreline from those probabilities and expected goals.
\end{quote}

\TODO{Verify the excerpted prompt wording against the final frozen prompt hashes and add the complete JSON schemas and special-question templates to the released appendix or repository.}

\end{document}
